\section{Certified computation and economic decisions}\label{sec:study49}
\subsection{The controlled investment family}
There are $d\in\{2,3,4\}$ capital coordinates on $[0,1]^d$, one investment action, and a common continuous shock. The feasible set is
\[
 0\le a\le\kappa(x)=\frac18+\frac1{8d}\sum_{j=1}^d x_j.
\]
With cyclic coordinate indices, $\gamma_j=1/2$ for odd $j$ and $1/4$ for even $j$, and $z\sim\mathrm{Unif}[-1/32,1/32]$, the dynamics are
\begin{equation}\label{eq:investment49}
 F_j(x,a,z)=\frac1{16}+\frac{x_j}2+\frac{x_{j+1}}8
       +\frac{x_j(1-x_{j+1})}{16}+\gamma_ja+(-1)^{j-1}z.
\end{equation}
Current costs are
\begin{equation}\label{eq:cost49}
 c(x,a)=\frac2d\sum_j(x_j-5/8)^2+
 \frac1{4d}\sum_j(x_j-x_{j+1})^2+
 2\left(\frac12-\frac2d\sum_jx_j\right)_+^2+pa^2+4a^4.
\end{equation}
Terminal cost doubles the first coefficient and retains the imbalance and shortage terms, with no action cost. Capital enters both adjustment incentives and feasible capacity. The nonlinear transition and shortage penalty couple the coordinates; the model is not separable linear-quadratic control. Prices $p\in\{1,4\}$ and horizons $T\in\{2,3\}$ are theoretical primitives, not calibrated estimates.

For the declared dimensions, valid primitive moduli in $\ell^1$ are
\[
 C_x=\frac{13}{2d},\quad L_g=\frac9d,\quad F_x=\frac{11}{16},\quad
 F_a=\sum_j\gamma_j,\quad C_a=\frac p2+\frac14,\quad K^A=\frac1{8d}.
\]
The transition is invariant by direct interval bounds. The state column sums of its Jacobian are at most $11/16$. The costs have piecewise-affine gradients; the cube vertices and the shortage-boundary vertices give the stated bounds. The supplemental proof restricts this explicit formula to the declared dimensions rather than presuming it for an arbitrary number of coordinates. The general constructive theorem uses whichever valid primitive moduli the economy supplies.

\subsection{Matched services and realized adaptation}
The execution-before protocol fixes three generators: compiled witness, uniform multilinear FVI, and curvature-refined coordinate FVI. Every generator forms its own future labels and receives the same state-node count, feasible action fractions and integration bins. For each two-state price--horizon cell the common resource ladder is
\[
 (N,K,M)=(4,2,1),(8,2,1),(16,4,2),(32,4,2),(64,8,4),(128,16,8).
\]
The entire ladder is run in three isolated processes per cell and method, rotating method order. This yields 36 two-state services and 216 rungs. State dimension three uses four declared rungs through $(32,8,4)$; dimension four uses three through $(16,8,4)$, with the witness and uniform generators, both horizons, price one and three repetitions. Those additional 24 services contain 84 rungs. The comparison across dimensions uses the same target five, not a dimension-specific certificate threshold. This is a finite common-accuracy experiment at a coarse tolerance, not evidence of an asymptotic or dimension-free rate.

The curvature comparator obtains an own-future pilot on a nine-node coordinate grid. For each axis it takes maxima of adjacent exact second differences across the other coordinates. It then bisects the interval maximizing squared length times the largest pilot curvature over that interval, with deterministic leftmost ties. A negligible positive floor permits a defined rule when the pilot is affine; the dominant unit baseline of the earlier heuristic is absent. Pilots, allocation and the realized axes are stored and charged. Realized nonuniformity, rather than the method's name, determines whether the experiment supplies an adaptive comparison. This remains a coordinate tensor algorithm; it is not presented as a sparse-grid or general adaptive-cell method.

Each rung starts from primitives. The internal complete-prefix clock includes own-future integration, all earlier failed rungs, pilots, exact compilation, actor construction, all-state certification, deployment checks and durable checkpoints. Separate process clocks include startup and the identical non-economic warm-up. A single CPU is selected, numerical-library threads are fixed to one, and CPU frequency is not controlled. Peak resident memory includes the process. Primitive query counts are component counts, not complete FLOPs or bit-operation counts. Independent policy evaluation has a separate joint service and is never added to an unrelated historical clock.

\input{tables/attainment49}
\input{tables/frontier49}
\input{tables/dimension49}
\input{results-discussion49}

\subsection{Actual first-crossing policies and replacement fees}
For the four two-state economic cells and each target $4,2,1$, the protocol selects repetition zero's first-crossing witness and uniform FVI policies when both attain. It also selects the particular favorable R48 target-two pair: $T=3,p=4$, witness at $N=32$ and FVI at $N=64$. Those old frozen policies are evaluated as an additional pair; their historical clocks are not reconstructed as part of the new service. Each selected pair is compared under the uniform initial law and point laws at $(1/8,1/8)$, $(1/2,1/2)$ and $(7/8,7/8)$, with full information or eight-bit observations. Finite observations use robust cell-wide capacity repair and downward action quantum $1/4096$.

Every sampled contrast uses 262,144 common interval paths, 40-bit equal initial and innovation bins, and outward mean and variance bounds. All new policy-pair and sensor contrasts share one family of at most 300 with error $1/100$ and log upper bound 13. The independent-bin model supplies the statistical assertion; the fixed PCG64 streams supply reproducibility. Old R48 intervals retain their separate original confidence event. Simultaneous interpretation of both old and new families is at least 98 percent by a union bound, not automatically 99 percent.

\input{tables/direct49}
The complete numerical supplement reports the directional replacement gains for every pair and the fee sufficient to certify that neither switch pays. It also reports the lower threshold that merely removes an identified profitable switch. The old $1/64$ decision is preserved, but no conclusion depends on treating that fee as calibrated. The deposited exact fee frontiers permit any nonnegative fee to be assessed on the same confidence event. Lower regret certificates are never described as observed policy-cost improvements.

\subsection{Purchasing information on the evaluated grid}
For each first-target-one controller, the protocol evaluates every integer precision $b=4,\ldots,16$ under the uniform initial law. Each policy is compared with its own 16-bit implementation; the reference self-comparison is exactly zero. Prices per coordinate-bit are $1/16384,1/4096,1/1024$. The deterministic all-state-account minimizer and the upper-net-cost selector of Theorem~\ref{thm:net49} use precisely that same grid. Neither is compared with an unobserved continuous optimum.

\input{tables/sensor49}
For each economic cell, the record additionally gives all nonnegative information prices by exact lower and upper affine envelopes. A retrospective price query within that family is allowed by simultaneous coverage and is not a new experiment. Initial-law dependence remains substantive: a uniform-state expected-cost decision need not be optimal at a particular initial state, and a conservatively optimal all-state bound may rationally purchase more information.

\subsection{The retained evidence and its interpretation}
The new service does not rewrite the original nonlinear training studies. The R42 tighter-target neural attainments, six of twelve scalar and five of six coupled services, remain capped-service outcomes; the R44 residual refinements of seven frozen failures remain diagnostics, not retraining successes. The 24 direct neural-minus-ridge intervals remain unresolved. The R45 constructive comparison still has nine of twelve neural and twelve of twelve spline attainments at its tightest target. Fixed64 and adaptive precision still certify all 36 historical precision services, with adaptive slower in 26 and about 0.98 percent slower in aggregate.

The R48 target-one count, twelve of twelve for witness versus six of twelve for each FVI method, also remains unchanged. At a relaxed tolerance 1.014 the old FVI near misses all attain; the release supplies a complete sensitivity table rather than reclassifying those original failures. The new common $N=128$ state rung uses separately declared action and integration budgets. It is not the unexecuted diagonal $(128,128,128)$ experiment. The original 48 constrained-policy intervals, including six favoring FVI and none favoring witness, remain available with newly exposed fee frontiers. Their adverse point-state pattern motivates inspection of actual policies; it is not removed by a favorable all-state certificate.

The complete development edition, unchanged historical companions and machine-readable preservation map retain all of these results and their proofs. This revision adds a controlled resource allocation and actual implementation decision, not a selectively favorable replacement record.
